\begin{table}[htbp]
    \caption{Firm-level heterogeneity in the debt-to-assets response}
    \label{tab:robustness_firm_heterogeneity}
    \centering
    \footnotesize
    \begin{adjustbox}{max width=\textwidth}
    \begin{tabular}{lrrr}
        \toprule
        Group & Effect (pp) & SE (pp) & \(p\)-value \\
        \midrule
        \multicolumn{4}{l}{\textit{Panel A: Debt-to-assets effects by pre-cutoff operational-risk quartile}} \\
        Q1, lowest risk & 0.31 & 1.25 & 0.808 \\
        Q2 & -1.33 & 0.96 & 0.182 \\
        Q3 & -1.70 & 1.65 & 0.315 \\
        Q4, highest risk & -3.06\(^{***}\) & 0.82 & 0.001 \\
        \addlinespace
        \multicolumn{4}{l}{\textit{Panel B: Debt-to-assets effects by pre-reform 30-percent accounting pressure}} \\
        Low pressure & -1.88\(^{**}\) & 0.81 & 0.031 \\
        High pressure & -1.46 & 0.96 & 0.145 \\
        \bottomrule
    \end{tabular}
    \end{adjustbox}
    \parbox{\textwidth}{\vspace{0.35em}\footnotesize\textit{Notes:} SE denotes standard error. The outcome is debt-to-assets. Estimates are in percentage points and come from pooled interaction models with the full leverage control set, firm and year fixed effects, and standard errors clustered by headquarters country. \(^{***}p<0.01\), \(^{**}p<0.05\), \(^{*}p<0.10\).}
\end{table}
